\documentclass[runningheads,a4paper]{llncs}
\usepackage[T1]{fontenc}
\usepackage[utf8]{inputenc}
\usepackage[portuguese]{babel}
\usepackage{lmodern}
\usepackage{booktabs}
\usepackage{tabularx}
\usepackage{graphicx}
\usepackage{tikz}
\usetikzlibrary{arrows.meta,fit,positioning}
\usepackage{pgfgantt}

\renewcommand{\keywordname}{{\bf Palavras-chave:}}
\renewcommand{\andname}{e}
\renewcommand{\lastandname}{\unskip\ e}
\setlength{\tabcolsep}{4pt}
\newcolumntype{L}[1]{>{\raggedright\arraybackslash}p{#1}}
\newcolumntype{Y}{>{\raggedright\arraybackslash}X}
\setlength{\textfloatsep}{8pt plus 2pt minus 2pt}
\setlength{\floatsep}{6pt plus 2pt minus 2pt}
\setlength{\intextsep}{6pt plus 2pt minus 2pt}

\begin{document}

\title{Implementação de mecanismos de QoS usando VNFs\\{\large Proposta de trabalho --- TP3}}
\titlerunning{QoS com VNFs --- Proposta (TP3)}

\author{André Miguel Rego Trindade Pinto \and Diogo Fernandes Pinto \and\\ Francisco José Fidalgo Lage}
\authorrunning{A. Pinto, D. Pinto e F. Lage}

\institute{Grupo PL11 --- Qualidade de Serviço na Internet,\\ Universidade do Minho, Braga, Portugal\\
\email{\{pg65400, pg65405, pg65407\}@alunos.uminho.pt}}

\maketitle

\begin{abstract}
Propõe-se um sistema de QoS do tipo DiffServ, construído como uma cadeia de quatro
funções de rede virtualizadas (VNFs) em contentores Docker sobre uma topologia emulada
semelhante à de um ISP, e avaliado com tráfego VoIP, de vídeo e de dados em massa.
\keywords{QoS \and NFV \and VNF \and DiffServ \and Cadeia de serviço}
\end{abstract}

\section{Objetivos e âmbito}

O objetivo é conceber, implementar e avaliar uma cadeia de serviço de QoS totalmente
virtualizada, na linha da arquitetura ETSI~NFV~\cite{etsi-nfv}. Cobrem-se os cinco
aspetos de QoS do enunciado, e não apenas os três exigidos: \emph{classificação},
\emph{controlo de acesso}, \emph{policiamento}, \emph{escalonamento} e
\emph{monitorização}. Para o sistema ser simples de explicar e de reproduzir, cada VNF
usa mecanismos nativos do \emph{kernel} Linux (\texttt{nftables} e \texttt{tc}) e todo
o ambiente é criado por \emph{scripts} no repositório Git do grupo
(\texttt{github.com/Daydream127/Qosi}).

\section{Topologia e cadeia de serviço}

A topologia (Fig.~\ref{fig:topo}) é emulada com Containernet~\cite{containernet}, uma
extensão do Mininet~\cite{mininet} que permite usar contentores Docker como nós, dentro
de uma máquina virtual Ubuntu (VirtualBox). Quatro clientes ligam-se a um
comutador Open~vSwitch (s1); todo o tráfego atravessa a cadeia de VNFs antes de chegar
ao encaminhador do núcleo (r1) e aos servidores. A ligação à saída da VNF3 é o
\emph{gargalo} (10~Mbit/s), onde a contenção entre classes se torna visível; na ligação
ao núcleo aplica-se \texttt{tc netem} para emular atraso, perdas e falhas.

\begin{figure}[!t]
\centering
\resizebox{\textwidth}{!}{%
\begin{tikzpicture}[font=\scriptsize, >=Latex,
  host/.style={draw, rounded corners=1pt, minimum width=1.3cm, minimum height=0.42cm, inner sep=1pt},
  sw/.style={draw, circle, inner sep=1.5pt, fill=gray!15},
  vnf/.style={draw, thick, fill=blue!8, minimum width=1.35cm, minimum height=1.0cm, align=center, inner sep=1pt},
  mgmt/.style={draw, fill=green!10, minimum height=0.5cm, inner sep=3pt, align=center}]
  % clientes
  \node[host] (hv) at (0,1.2) {h\_voip};
  \node[host] (hd) at (0,0.4) {h\_video};
  \node[host] (hb) at (0,-0.4) {h\_bulk};
  \node[host, densely dashed] (hx) at (0,-1.2) {h\_x};
  \node[above=1pt of hv, font=\scriptsize\itshape] {Clientes};
  \node[sw] (s1) at (1.4,0) {s1};
  % cadeia
  \node[vnf] (c) at (2.9,0) {VNF1\\Classificador};
  \node[vnf] (f) at (4.6,0) {VNF2\\Firewall};
  \node[vnf] (p) at (6.3,0) {VNF3\\Policer/\\Escalonador};
  \node[vnf] (m) at (9.3,0) {VNF4\\Monitor};
  \node[draw, dashed, inner sep=5pt, fit=(c)(f)(p)(m),
        label={[font=\scriptsize\itshape]above:Cadeia de serviço (contentores Docker)}] (chain) {};
  % núcleo e servidores
  \node[sw] (r1) at (11.6,0) {r1};
  \node[sw] (s2) at (12.5,0) {s2};
  \node[host] (sv) at (13.85,0.8) {srv\_voip};
  \node[host] (sd) at (13.85,0) {srv\_video};
  \node[host] (sb) at (13.85,-0.8) {srv\_bulk};
  \node[above=1pt of sv, font=\scriptsize\itshape] {Servidores};
  % ligações de dados
  \foreach \h in {hv,hd,hb,hx} \draw (\h.east) -- (s1);
  \draw[->] (s1) -- (c);
  \draw[->] (c) -- (f);
  \draw[->] (f) -- (p);
  \draw[->, very thick] (p) -- node[above] {10 Mbit/s} node[below] {gargalo} (m);
  \draw[->] (m) -- node[above, pos=0.66] {netem} (r1);
  \draw (r1) -- (s2);
  \foreach \s in {sv,sd,sb} \draw (s2) -- (\s.west);
  % plano de gestão
  \node[mgmt] (prom) at (6.0,-1.75) {Prometheus + Grafana (rede de gestão)};
  \foreach \v in {c,f,p,m} \draw[dashed, ->, gray] (\v.south) -- (\v.south |- prom.north);
\end{tikzpicture}}
\caption{Topologia emulada e cadeia de serviço (h\_x é uma origem não autorizada).}
\label{fig:topo}
\end{figure}

As VNFs são encadeadas em linha: cada contentor tem duas interfaces de dados e
funciona como um salto IP, com encaminhamento estático definido pelo \emph{script} de
topologia, que assume o papel de orquestrador. A ordem da cadeia e os parâmetros de
cada VNF (Tabela~\ref{tab:vnfs}) ficam num ficheiro de configuração.

\begin{table}[!ht]
\caption{VNFs da cadeia, aspetos de QoS cobertos e ferramentas.}
\label{tab:vnfs}
\centering
\footnotesize
\begin{tabularx}{\textwidth}{@{}L{3.5cm} L{2.4cm} Y@{}}
\toprule
VNF e \emph{aspeto de QoS} & Ferramenta & Função \\
\midrule
1 Classificador\newline \emph{Classificação} & \texttt{nftables} &
  Identifica fluxos por endereços, protocolo e portas; marca o DSCP. \\
2 Firewall\newline \emph{Controlo de acesso} & \texttt{nftables} &
  Política \emph{default-deny}: só admite origens e serviços autorizados. \\
3 Policer/Escalonador\newline \emph{Policiamento, escalonamento} & \texttt{tc} (police, HTB) &
  \emph{Token bucket} por classe; prioridade para EF e partilha ponderada entre AF31
  e BE. \\
4 Monitor\newline \emph{Monitorização} & Python, Prometheus, Grafana &
  Contadores por classe (passivo) e sondas com n.º de sequência e estampilha temporal
  (ativo). \\
\bottomrule
\end{tabularx}
\end{table}

\section{Cenários de tráfego}

O tráfego é sintético, gerado com \texttt{iperf3} e com um gerador de sondas próprio em
Python. As classes e a marcação DSCP seguem o DiffServ~\cite{rfc2475} e o
RFC~4594~\cite{rfc4594}; os objetivos de VoIP são os da classe~0 da
Y.1541~\cite{y1541} (Tabela~\ref{tab:trafego}).

\begin{table}[!ht]
\caption{Classes de tráfego, perfis gerados e objetivos de QoS.}
\label{tab:trafego}
\centering
\footnotesize
\begin{tabularx}{\textwidth}{@{}l Y l Y@{}}
\toprule
Aplicação & Perfil gerado & DSCP & Objetivo \\
\midrule
VoIP (G.711) & UDP, 50~pacotes/s de 200~B ($\approx$80~kbit/s) & EF (46) &
  IPTD $\le$ 100~ms, IPDV $\le$ 50~ms, IPLR $\le 10^{-3}$ \\
Vídeo & \emph{Streaming} UDP de débito constante, 4~Mbit/s & AF31 (26) &
  Débito sustentado, IPLR $\le 10^{-3}$ \\
Dados em massa & TCP, 1 a $N$ fluxos sem limite de débito & BE (0) &
  Sem garantias; usa a capacidade restante \\
Não autorizado & UDP a partir de h\_x & --- & Bloqueado na VNF2 \\
\bottomrule
\end{tabularx}
\end{table}

\section{Metodologia de avaliação}

Por classe, medem-se as métricas \emph{one-way} estudadas na UC: atraso (OWD ou IPTD),
variação do atraso (IPDV), perdas (IPLR) e débito. Como todos os nós emulados partilham
o relógio da máquina, o atraso num só sentido é medido sem erro de sincronização. Para
VoIP estima-se ainda o MOS a partir do atraso e das perdas, com o modelo~E
simplificado~\cite{cole}.

Comparam-se quatro configurações: (i)~sem QoS, com as VNFs em modo de passagem;
(ii)~classificação e escalonamento; (iii)~com policiamento; (iv)~cadeia completa, com
tráfego não autorizado. Em cada uma varia-se a carga (número de fluxos TCP e débito de
vídeo) e as condições da ligação ao núcleo (atraso, perdas e falha temporária). Cada
cenário é repetido cinco vezes durante 60~s, reportando-se médias com intervalos de
confiança de 95\%. Um único comando lança a campanha e produz os CSV e os gráficos. Na
otimização ajustam-se os \emph{token buckets} e os pesos das filas, documentando os
compromissos e a sobrecarga da cadeia.

\section{Planeamento}

O plano (Fig.~\ref{fig:gantt}) segue as quatro fases e os três marcos do enunciado.

\begin{figure}[!ht]
\centering
\begin{ganttchart}[
  x unit=0.5cm, y unit chart=0.3cm, y unit title=0.4cm,
  title height=1, bar height=0.55, group height=0.25, group peaks height=0.15,
  title label font=\scriptsize\bfseries,
  bar label font=\scriptsize,
  group label font=\scriptsize\bfseries,
  group label node/.append style={xshift=-3pt},
  milestone label font=\scriptsize\itshape,
  bar/.append style={fill=blue!25},
  milestone/.append style={fill=red!70},
  vgrid={*1{draw=none}, *1{dotted}}
]{1}{14}
\gantttitle{S1}{2}\gantttitle{S2}{2}\gantttitle{S3}{2}\gantttitle{S4}{2}%
\gantttitle{S5}{2}\gantttitle{S6}{2}\gantttitle{S7}{2}\\
\ganttgroup{Fase 1: ambiente e plano}{1}{6}\\
\ganttbar{Estudo do enunciado e das ferramentas}{1}{4}\\
\ganttbar{Arquitetura, VM e topologia base}{5}{6}\\
\ganttmilestone{M1: proposta (13 de novembro)}{6}\\
\ganttgroup{Fase 2: VNFs isoladas}{7}{9}\\
\ganttbar{VNF1 Classificador e VNF2 Firewall}{7}{7}\\
\ganttbar{VNF3 Policer/Escalonador}{7}{8}\\
\ganttbar{VNF4 Monitor e geradores de tráfego}{8}{9}\\
\ganttgroup{Fase 3: integração e otimização}{9}{12}\\
\ganttbar{Integração da cadeia e primeiros testes}{9}{10}\\
\ganttmilestone{M2: relatório de progresso (fim de S5)}{10}\\
\ganttbar{Congestionamento, falhas e otimização}{11}{12}\\
\ganttgroup{Fase 4: análise e relatório}{12}{14}\\
\ganttbar{Campanha de experiências e gráficos}{12}{13}\\
\ganttbar{Relatório final e demonstração}{13}{14}\\
\ganttmilestone{M3: entrega final (fim de S7)}{14}
\end{ganttchart}
\caption{Planeamento por semanas (S1 a S7) e marcos de entrega.}
\label{fig:gantt}
\end{figure}

\begin{thebibliography}{8}
\bibitem{etsi-nfv} ETSI: Network Functions Virtualisation (NFV); Architectural
Framework. ETSI GS NFV 002 V1.2.1 (2014)
\bibitem{containernet} Peuster, M., Karl, H., van Rossem, S.: MeDICINE: Rapid prototyping
of production-ready network services in multi-PoP environments. In: IEEE NFV-SDN (2016)
\bibitem{mininet} Lantz, B., Heller, B., McKeown, N.: A network in a laptop: rapid
prototyping for software-defined networks. In: ACM HotNets-IX (2010)
\bibitem{rfc2475} Blake, S., et al.: An Architecture for Differentiated Services.
RFC 2475, IETF (1998)
\bibitem{rfc4594} Babiarz, J., Chan, K., Baker, F.: Configuration Guidelines for DiffServ
Service Classes. RFC 4594, IETF (2006)
\bibitem{y1541} ITU-T: Network performance objectives for IP-based services.
Recommendation Y.1541 (2011)
\bibitem{cole} Cole, R.G., Rosenbluth, J.H.: Voice over IP performance monitoring.
ACM SIGCOMM Computer Communication Review \textbf{31}(2), 9--24 (2001)
\end{thebibliography}

\end{document}
